
\chapter{Review of Related Literature and Studies}

\label{chap:review}

This chapter presents the concepts, theories, literature, and previous studies relevant to the development of a Shallow Water Equation-based Physics-Informed Neural Network for flood prediction. Since the study combines several fields, namely mathematics, physics, machine learning, and flood prediction, the discussion is organized to establish the theoretical and scientific foundation of the proposed model.

\section{Preliminaries}

This section introduces the fundamental concepts necessary for understanding the study. It provides an overview of flooding, hydrodynamic modeling, artificial neural networks, and Physics-Informed Neural Networks.

\subsection{Fundamental Concepts of Flooding}

Discuss the basic concepts of flooding, including its causes, characteristics, and effects. Explain important flood-related variables such as water depth, flow velocity, rainfall, and surface runoff.

\subsection{Hydrodynamic Modeling}

Discuss hydrodynamic modeling and its role in simulating the movement and distribution of water. Introduce the use of mathematical equations and numerical methods in representing flood behavior.

\subsection{Artificial Neural Networks}

Introduce artificial neural networks, including their basic structure, components, and learning process. Explain how neural networks can be used to approximate complex mathematical relationships.

\subsection{Physics-Informed Neural Networks}

Discuss Physics-Informed Neural Networks and how physical laws and governing equations are incorporated into neural network training. Explain their advantages and limitations compared with purely data-driven approaches.

\section{Theoretical Background}

This section presents the mathematical and computational theories that serve as the foundation of the proposed SWE-based PINN.

\subsection{Shallow Water Equations}

Present the Shallow Water Equations and explain their application in modeling flood flow and water movement over terrain.

\subsection{Conservation of Mass}

Discuss the principle of conservation of mass and its mathematical representation in the Shallow Water Equations.

\subsection{Conservation of Momentum}

Explain the principle of conservation of momentum and its role in describing the movement and velocity of water.

\subsection{Governing Equations for Flood Flow}

Present the governing equations used to describe flood flow, including the relevant flux terms and source terms associated with terrain slope and friction.

\subsection{Physics-Informed Neural Network Framework}

Discuss the mathematical structure of the PINN, including its neural network architecture, input and output variables, automatic differentiation, and incorporation of the governing equations.

\subsection{Loss Function of a PINN}

Explain the components of the PINN loss function, including the data loss, initial and boundary condition losses, and physics-based residual loss. Discuss how these components guide the network toward physically consistent solutions.

\section{Review of Related Literature}

This section discusses existing literature, concepts, and theoretical discussions related to flood prediction, hydrodynamic modeling, and the application of machine learning and Physics-Informed Neural Networks.

\subsection{Flood Prediction and Modeling}

Review the existing approaches used in flood prediction and modeling, including traditional mathematical and computational techniques.

\subsection{Flood Modeling in Urban Areas}

Discuss the challenges of urban flood modeling, including complex terrain, land use, drainage systems, and surface runoff.

\subsection{Shallow Water Equation Applications}

Review the applications of the Shallow Water Equations in simulating water movement, flood propagation, and water depth distribution.

\subsection{Physics-Informed Neural Networks in Hydrology}

Discuss the use of PINNs in hydrological modeling, particularly their ability to incorporate physical constraints into data-driven prediction models.

\subsection{Machine Learning for Flood Prediction}

Review machine learning approaches used in flood prediction, including artificial neural networks and other data-driven methods. Discuss their strengths and limitations.

\subsection{Flood Risk in Pasig City and Metro Manila}

Discuss the flood-related conditions and challenges in Pasig City and Metro Manila. Include relevant information about rainfall, river systems, urban development, and flood susceptibility.

\section{Review of Related Studies}

This section presents previous research studies that are directly related to the proposed research. Focus on the methods, datasets, results, and findings of each study.

\subsection{SWE-Based Flood Prediction Studies}

Review previous studies that used the Shallow Water Equations for flood simulation and prediction. Discuss their numerical methods, model performance, and findings.

\subsection{PINN-Based Flood Prediction Studies}

Present research studies that implemented Physics-Informed Neural Networks for flood prediction. Discuss their model architectures, governing equations, training methods, and evaluation metrics.

\subsection{PINN Applications in Hydrodynamics}

Review studies that applied PINNs to hydrodynamic problems, particularly those involving fluid flow, water movement, and conservation laws.

\subsection{Flood Prediction Studies in the Philippines}

Present previous research conducted in the Philippines concerning flood prediction, flood modeling, and flood risk assessment. Discuss the methodologies and findings relevant to the selected study area.

\section{Synthesis of Literature}

Summarize the common findings, approaches, and limitations identified in the reviewed literature and studies. Discuss the similarities and differences between traditional hydrodynamic models, machine learning approaches, and Physics-Informed Neural Networks.

\section{Research Gap}

Identify the limitations and unresolved issues in previous research. Explain the need for further investigation into the convergence, approximation error, and physical consistency of SWE-based Physics-Informed Neural Networks for flood modeling.

\section{Relevance to the Present Study}

Explain how the reviewed literature and studies support the development of the proposed SWE-based PINN. Establish the connection between the existing research and the objectives of the present study, particularly its comparison with a classical numerical method.

\section{Chapter Summary}

Provide a summary of the important concepts, theories, literature, and studies discussed in this chapter. Highlight how these establish the foundation for the methodology and model development presented in the succeeding chapter.
